%% kapa: exam=1 round=33 year=2022 period=1 form=A subject=economics subject_ko=경제학원론 seq=24 number=64
%% source: Q-Net 감정평가사 기출문제 (cst00309) artlSeq=5209997 file=raw/1008/5209997/2022년도 제33회 감정평가사 1차시험 1교시(민법,경제학,부동산학).pdf
%% conversion: kordoc 4.24.51 -> tools/gen1.py (auto); review=proofread-sonnet; flags=none
\begin{question}{64}
A국의 단기 필립스 곡선이 아래와 같을 때 이에 관한 설명으로 옳지 않은 것은? (단, $\pi$, $\pi^e$, $u$, $u_n$은 각각 인플레이션율, 기대 인플레이션율, 실업률, 자연 실업률이다.)
\begin{framedtext}
$\pi - \pi^e = -0.5(u - u_n)$
\end{framedtext}
\begin{choices}
  \item 총공급곡선이 수직선인 경우에 나타날 수 있는 관계이다.
  \item 총수요 충격이 발생하는 경우에 나타날 수 있는 관계이다.
  \item 인플레이션율과 실업률 사이에 단기적으로 상충관계가 있음을 나타낸다.
  \item 고용이 완전고용수준보다 높은 경우에 인플레이션율은 기대 인플레이션율보다 높다.
  \item 인플레이션율을 1\%p 낮추려면 실업률은 2\%p 증가되어야 한다.
\end{choices}
\end{question}
\answer{1}
